\documentclass[12pt,a4paper]{ctexart}
\usepackage[margin=1in]{geometry}
\usepackage{graphicx}
\usepackage[dvipsnames]{xcolor}
\usepackage{amsmath}
\usepackage{amssymb}
\usepackage{booktabs}
\usepackage{array}
\usepackage{enumitem}
\usepackage{listings}
\usepackage{hyperref}

\definecolor{codebg}{HTML}{F7F8FA}
\definecolor{codeframe}{HTML}{D9DEE7}
\definecolor{codekeyword}{HTML}{0066CC}
\definecolor{codestring}{HTML}{A31515}
\definecolor{codecomment}{HTML}{008000}

\lstdefinestyle{reportcode}{
  backgroundcolor=\color{codebg},
  basicstyle=\ttfamily\small,
  frame=single,
  rulecolor=\color{codeframe},
  frameround=ffff,
  columns=fullflexible,
  breaklines=true,
  breakatwhitespace=true,
  tabsize=4,
  showstringspaces=false,
  keepspaces=true,
  keywordstyle=\color{codekeyword}\bfseries,
  commentstyle=\color{codecomment},
  stringstyle=\color{codestring},
  numbers=left,
  numberstyle=\tiny\color{gray},
  numbersep=8pt,
  xleftmargin=1.8em,
  framexleftmargin=1.2em,
  aboveskip=0.8em,
  belowskip=0.8em
}

\lstset{style=reportcode}
\hypersetup{
  colorlinks=true,
  linkcolor=MidnightBlue,
  urlcolor=RoyalBlue,
  citecolor=MidnightBlue
}

\title{并行程序设计与算法实验\\实验四：基于 OpenMP 的并行计算}
\author{姓名：元朗曦 \quad 学号：23336294}
\date{\today}

\begin{document}
\maketitle

\section{实验目的}
本实验的目标是使用 OpenMP 实现基于共享内存的并行计算，包括矩阵乘法和数组求和两个典型算法。通过本实验希望达到以下目的：
\begin{enumerate}[itemsep=0.3em]
  \item 掌握 OpenMP 的基本编程模型和指令集，特别是 \texttt{\#pragma omp parallel for} 的使用。
  \item 理解并行循环中的工作分配、负载均衡与 \texttt{schedule} 策略的影响。
  \item 学习规约（reduction）操作在数据并行中的应用，避免竞争条件。
  \item 通过性能测试验证线程数增加对程序运行时间的影响，分析加速比和并行效率。
  \item 对比 OpenMP 与传统线程库（如 POSIX 线程）的编程便利性。
\end{enumerate}

\section{实验原理}

\subsection{OpenMP 并行编程模型}
OpenMP（Open Multi-Processing）是一个基于共享内存的并行编程框架，通过编译指导语句（pragma）和运行时库函数，提供高层次的并行抽象。相比 POSIX 线程库（pthread），OpenMP 提供了更简洁的 API，程序员无需手动创建和销毁线程，编译器会自动生成并行代码。

OpenMP 的核心编程模型基于"fork-join"范式：主线程执行串行代码，在遇到并行指令时创建工作线程，所有线程并行执行指定的代码块，然后在隐含屏障（implicit barrier）处同步，主线程继续执行后续串行代码。

\subsection{并行循环与工作分配}
在数据并行中，最常见的模式是使用 \texttt{\#pragma omp parallel for} 将循环迭代分配给多个线程。调度策略通过 \texttt{schedule()} 子句指定：

\begin{enumerate}[itemsep=0.2em]
  \item \textbf{静态调度}（schedule(static)）：将迭代均匀分配给各线程，无运行时开销，但若迭代工作量不均则负载不均衡。
  \item \textbf{动态调度}（schedule(dynamic)）：运行时将迭代动态分配，负载均衡好但同步开销大。
  \item \textbf{自适应调度}（schedule(guided)）：介于两者之间，初期大块分配，后期逐步减小块大小。
\end{enumerate}

对于矩阵乘法中行与行之间独立的外层循环，使用静态调度足以获得良好的负载均衡。

\subsection{规约操作}
数组求和等操作涉及多个线程的累积。直接在并行区域中共享变量会产生竞争条件。OpenMP 提供 \texttt{reduction(op:var)} 子句来安全地累积结果：编译器为每个线程创建临时副本，各线程独立计算，最后通过指定的操作（如 \texttt{+}）合并结果。

\subsection{浮点性能指标}
浮点计算性能用 GFLOPS（十亿浮点运算次数/秒）表示：

\begin{equation}
\text{GFLOPS} = \frac{\text{浮点运算次数}}{t \times 10^9}
\end{equation}

对于矩阵乘法 $C = A \times B$，其中 $A \in \mathbb{R}^{m \times n}$，$B \in \mathbb{R}^{n \times k}$，浮点运算次数约为 $2mnk$（$n$ 次乘-加操作，每个迭代）。

加速比定义为：

\begin{equation}
\text{Speedup} = \frac{t_{\text{serial}}}{t_{\text{parallel}}}
\end{equation}

理想情况下，用 $p$ 个线程运行程序的加速比应接近 $p$（线性加速），但实际受到通信开销、负载不均衡、Amdahl's Law 等限制。

\section{实验内容}

\subsection{程序设计}
本实验在 \texttt{lab/04/src} 中实现两个核心程序：

\subsubsection{矩阵乘法 (matmul\_openmp)}
\begin{itemize}
  \item 功能：计算 $C = A \times B$，其中 $A \in \mathbb{R}^{m \times n}$，$B \in \mathbb{R}^{n \times k}$。
  \item 输入：命令行参数 \texttt{<m> <n> <k> <threads> [--verify]}。
  \item 输出：矩阵形状、线程数、运行时间、GFLOPS、可选的校验结果。
  \item 并行策略：使用 \texttt{\#pragma omp parallel for schedule(static)} 并行化外层行迭代。
  \item 校验：与串行版本比较，检查数值精度（最大绝对误差 < $10^{-8}$）。
\end{itemize}

\subsubsection{数组求和 (array\_sum\_openmp)}
\begin{itemize}
  \item 功能：计算数组元素和 $\sum_{i=0}^{N-1} a[i]$。
  \item 输入：命令行参数 \texttt{<size> <threads> [--verify]}。
  \item 输出：数组规模、线程数、运行时间、吞吐量（Gelem/sec）、结果、可选的校验。
  \item 并行策略：使用 \texttt{\#pragma omp parallel for reduction(+:sum)} 安全地累积和。
  \item 校验：与串行版本比较，检查精度（误差 < $10^{-6}$）。
\end{itemize}

\subsection{构建与运行}
使用 CMake 作为构建工具：

\begin{lstlisting}[language=bash]
# 配置和构建
mkdir -p lab/04/build
cd lab/04/build
cmake ..
cmake --build . -j

# 运行矩阵乘法（512x512x512, 4线程, 含验证）
./src/matmul_openmp 512 512 512 4 --verify

# 运行数组求和（1000万元素, 4线程, 含验证）
./src/array_sum_openmp 10000000 4 --verify

# 性能测试：不同线程数下的加速比
for t in 1 2 4 8; do
  ./src/matmul_openmp 1024 1024 1024 $t
done
\end{lstlisting}

\section{实验结果}

\subsection{正确性验证}
两个程序都通过了数值校验测试，验证无误。

\subsubsection{矩阵乘法校验 (512×512×512, 4线程)}
\begin{lstlisting}
Task: matrix_multiplication_openmp
Shape: A(512x512), B(512x512)
Threads: 4
Time (sec): 0.162897
GFLOPS: 1.648
Verify: PASS, max_abs_diff=0.0000000000
\end{lstlisting}

所有元素与串行版本完全相同（误差在浮点精度范围内）。

\subsubsection{数组求和校验 (1000万元素, 4线程)}
\begin{lstlisting}
Task: array_sum_openmp
Size: 10000000
Threads: 4
Time (sec): 0.006282
Throughput (Gelem/sec): 1.591899
Result: 1789.129084
Verify: PASS, serial_result=1789.1290836077, diff=0.0000000001
\end{lstlisting}

规约操作正确，避免了竞争条件。

\subsection{性能分析：矩阵乘法加速比}
测试环境：$1024 \times 1024 \times 1024$ 矩阵乘法，不同线程数下的性能表现。

\begin{table}[h]
\centering
\caption{OpenMP 矩阵乘法性能扩展性}
\begin{tabular}{cccccc}
\toprule
\textbf{线程数} & \textbf{时间(秒)} & \textbf{GFLOPS} & \textbf{加速比} & \textbf{效率(\%)} \\
\midrule
1  & 4.950  & 0.434  & 1.00  & 100.0 \\
2  & 2.518  & 0.853  & 1.97  & 98.3  \\
4  & 1.298  & 1.655  & 3.81  & 95.3  \\
8  & 0.860  & 2.496  & 5.75  & 71.9  \\
\bottomrule
\end{tabular}
\end{table}

\subsection{结论与分析}

\subsubsection{线性加速的实现}
从 1 线程到 4 线程，加速比接近线性（1×、2×、4×），平行效率分别为 100\%、98.3\%、95.3\%，说明在这个范围内，OpenMP 的开销很小，并行收益显著。

\subsubsection{8 线程处的性能下降}
8 线程时加速比为 5.75×（而非理想的 8×），平行效率降至 71.9\%。这是因为：
\begin{enumerate}[itemsep=0.2em]
  \item 当前测试环境可能只有 4 个物理核心；过度订阅（oversubscription）会增加上下文切换开销。
  \item OpenMP 线程创建、销毁及同步的开销在线程数增多时变得相对更显著。
  \item 缓存行竞争（cache line contention）在高线程数下加剧。
\end{enumerate}

\subsubsection{与 pthread 的对比}
相比 lab/03 的 pthread 实现（matmul\_pthreads 最多 16 线程），OpenMP 提供了以下优势：
\begin{itemize}
  \item \textbf{编程简洁性}：无需手动管理线程生命周期，只需一条 pragma 指令。
  \item \textbf{灵活的调度策略}：通过 \texttt{schedule} 子句可轻易调整工作分配方式。
  \item \textbf{规约支持}：内置 \texttt{reduction} 机制，避免竞争条件和锁竞争。
  \item \textbf{可移植性}：基本 OpenMP 功能在各种编译器上都有支持。
\end{itemize}

同时，OpenMP 对共享内存架构的优化比消息传递（MPI）更直接，适合单机多核、SMP 系统。

\subsubsection{实验收获}
\begin{enumerate}[itemsep=0.3em]
  \item 掌握了 OpenMP 的核心用法，能够快速将串行程序并行化。
  \item 理解了共享内存编程中规约、同步、负载均衡的关键概念。
  \item 通过实测加速比，验证了理论上的加速上界与实际性能的差异。
  \item 认识到无论多核还是众核系统，仍需关注缓存、负载均衡、同步开销等细节。
\end{enumerate}

\subsection*{总结}
本实验成功演示了 OpenMP 在矩阵乘法和数组求和中的应用。在线程数不超过物理核心数时，程序表现出接近线性的加速比；超出物理核心数后，加速比增长放缓。OpenMP 的高级抽象大幅降低了并行编程难度，同时保持了足够的性能，是现代共享内存并行计算的主流选择。

\end{document}
